\chapter{Fundamental, Analyst and Event Features}\label{ch:rs:fundamental-analyst-and-event-features}

A company closes its quarter on 31 December, announces its earnings on 25 January and files its quarterly report on 5
February. A database keyed on the period's end shows the quarter's earnings on 31 December, and a backtest built on it trades
the announcement's news three and a half weeks before anyone could read it. The error is invisible in the code, which joins
on a date column like any other, and spectacular in the results. This chapter keys accounting numbers, surprises, analyst
forecasts and event dates on the time they became known. Its real data are ten large US companies' reported earnings and
filing calendars from the SEC's public EDGAR interfaces; its measurement of the damage is on \texttt{firm.synthmkt}, where the
truth is known.

\section{Accounting data, point-in-time}

\begin{definition}[Fiscal period, filing lag]\label{def:rs:fundamental-analyst-and-event-features:fiscal}
A \emph{fiscal period}\index{fiscal period} is the quarter or year an accounting report describes, identified by its start and
end dates; the \omterm{def:rs:point-in-time-data-and-the-biases:bitemporal}{valid time} of every number in the report is the period. The \emph{filing lag}\index{filing lag} is the time from
the period's end to the filing of the report with the regulator.
\end{definition}

An accounting number has at least three dates after its period ends: the earnings release (a press release, filed in the US as a
Form 8-K under item 2.02), the periodic report (Form 10-Q for a quarter, 10-K for a year), and any later filing that reports the
same period again as a comparative, possibly with a different value
(\cref{fig:rs:fundamental-analyst-and-event-features:timeline}). The market learns the number at the release; a database built
from the periodic reports learns it at the filing; a database that stores only the period end pretends to know it at the end.

\begin{omfigure}
\begin{tikzpicture}[font=\small,x=0.19cm]
\draw[thick,-Stealth] (0,0) -- (66,0) node[right]{days};
\foreach \x/\l in {0/0,23/23,29/29,40/40,45/45,60/60} {\draw (\x,0.12) -- (\x,-0.12) node[below]{\l};}
\node[above,align=center,text=omThm] at (0,0.35) {period\\ends};
\node[above,align=center,text=omDef] at (23,0.35) {earnings\\release};
\node[above,align=center,text=omDef] at (31,1.35) {10-Q\\filed};
\draw[omDef] (29,0.15) -- (30,1.3);
\draw[black!50,dashed] (40,-0.6) -- (40,2.3) node[above,align=center,font=\footnotesize,text=black!70]{10-Q deadline,\\large filers};
\draw[black!50,dashed] (45,-0.6) -- (45,1.2) node[above,align=center,font=\footnotesize,text=black!70]{others};
\draw[black!50,dashed] (60,-0.6) -- (60,1.2) node[above,align=center,font=\footnotesize,text=black!70]{10-K,\\large filers};
\draw[omThm,thick,Stealth-Stealth] (0,-1.1) -- node[below,font=\footnotesize]{head start of a period-end key} (23,-1.1);
\end{tikzpicture}

\omcaption{The dates of a quarterly number, at the medians of ten large US companies' first three fiscal quarters since 2013 (23 days to
the earnings release, 29 to the 10-Q), with the SEC's filing deadlines.}\label{fig:rs:fundamental-analyst-and-event-features:timeline}
\end{omfigure}

\begin{dated}{2026-09}{US periodic report deadlines}\label{dat:rs:fundamental-analyst-and-event-features:deadlines}
The SEC's Form 10-Q instructions require the quarterly report within 40 days of the quarter's end for large accelerated and accelerated
filers and 45 days for all other registrants; no 10-Q is filed for the fourth quarter. Form 10-K is due 60 days after the fiscal year's end
for large accelerated filers, 75 for accelerated filers and 90 for all others. The SEC's EDGAR application programming interfaces serve
every filer's submission history and the XBRL data of its financial statements as JSON, updated in real time, without authentication or keys.
\end{dated}

The ten companies of \texttt{rs\_fetch\_edgar.py} (Apple, Microsoft, Walmart, Coca-Cola, Johnson \& Johnson, Exxon Mobil,
Procter \& Gamble, Intel, Home Depot and Nike) have 542 \omterm{def:rs:fundamental-analyst-and-event-features:fiscal}{fiscal period} ends since 2013 in their EDGAR submission histories, 541 of
them followed by an earnings release within 90 days. For the first three quarters the median company released its earnings 23 days
after the period ended and filed its 10-Q 29 days after; for fiscal years the medians are 26 and 48 days. The largest 10-Q lag is
41 days and the largest 10-K lag 60, inside the deadlines. Only 16\% of reports were filed on the day of the release; the median gap is
a week, during which the number is public but absent from a database fed by filings
(\cref{fig:rs:fundamental-analyst-and-event-features:calendar}).

\begin{omfigure}
\begin{tikzpicture}
\begin{axis}[width=11cm,height=6.2cm,xbar,bar width=5pt,xlabel={median days after the quarter's end},
  tick label style={font=\small},label style={font=\small},xmin=0,xmax=42,ymin=-0.6,ymax=9.6,ytick={0,1,2,3,4,5,6,7,8,9},
  yticklabels={Home Depot,Johnson \& Johnson,Walmart,Nike,Procter \& Gamble,Coca-Cola,Intel,Microsoft,Exxon Mobil,Apple},
  y dir=reverse,grid=major,grid style={gray!25},legend columns=2,
  legend style={font=\footnotesize,at={(0.5,-0.2)},anchor=north,draw=none,fill=none,
  /tikz/every even column/.append style={column sep=8pt}},table/col sep=comma]
\addplot[fill=omDef!55,draw=omDef] table[x=release,y=k]{figdata/research/11-fundamental-analyst-and-event-features/calendar.csv};
\addplot[fill=gray!35,draw=gray] table[x=filing,y=k]{figdata/research/11-fundamental-analyst-and-event-features/calendar.csv};
\draw[black!60,dashed] (axis cs:40,-0.6) -- (axis cs:40,9.6);
\legend{earnings release (8-K 2.02),10-Q filed}
\end{axis}
\end{tikzpicture}

\omcaption{Median days from the end of each of the first three fiscal quarters to the earnings release and to the 10-Q filing, ten
companies, 2013 to 2026; the dashed line is the 40-day deadline. Some file with the release, some weeks later. Data: SEC EDGAR
submissions, via \texttt{rs\_fetch\_edgar.py}.}\label{fig:rs:fundamental-analyst-and-event-features:calendar}
\end{omfigure}

The numbers also change after they are first filed. Each XBRL fact in EDGAR lists every filing that reported it, so a quarter's
diluted earnings per share can be followed from its first report through every later comparative. Among the ten companies, 31
quarterly values changed:

\begin{center}
\small
\begin{tabular}{@{}lccl@{}}
\toprule
company & values changed & by an exact ratio & ratios\\
\midrule
Apple & 13 & 13 & 7 and 4\\
Nike & 7 & 6 & 2\\
Microsoft & 5 & 0 & ---\\
Walmart & 3 & 3 & 3\\
Coca-Cola & 2 & 2 & 2\\
Johnson \& Johnson & 1 & 0 & ---\\
the other four & 0 & 0 & ---\\
\bottomrule
\end{tabular}
\end{center}

Twenty-four of the 31 changed by an exact ratio (to the cent), the signature of a share split: Apple's 10-Ks record a seven-for-one split
on 6 June 2014 and a four-for-one split on 28 August 2020, after which all per-share information was retroactively adjusted. The other
seven are \omterm{def:rs:point-in-time-data-and-the-biases:vintage}{restatements} of substance. Microsoft's diluted earnings for its fiscal year 2017 were first reported as \$2.71 and a year
later as \$3.25: it adopted a new revenue standard in its fiscal 2018 using the full retrospective method, which required it to restate
each prior period presented. A backtest of fiscal 2017 must use \$2.71 until August 2018.

\begin{definition}[Trailing twelve months]\label{def:rs:fundamental-analyst-and-event-features:ttm}
The \emph{trailing twelve months}\index{trailing twelve months} (TTM) value of a flow (earnings, revenue, cash flow) at a quarter's end is
the sum of the four most recent consecutive quarters, each as known at the decision time; a fourth quarter that is reported only inside
an annual figure is the year minus the first nine months.
\end{definition}

Summing four quarters is where the dates bite together. Apple's first-reported quarterly diluted EPS for its fiscal 2014 were \$14.50 and
\$11.62 (before the June split) and \$1.28 and \$1.42 (after): a TTM built from first reports is \$28.82, while the annual figure Apple
reported for the same year was \$6.45. The same happens across the 2020 split, \$10.85 against an annual \$3.28
(\cref{fig:rs:fundamental-analyst-and-event-features:ttm}). Today's view has the opposite problem: it rescales history with knowledge
of later splits, and it is itself a patchwork, because a period is rescaled only if a later filing reported it again. The rule is to
store per-share values as first reported with their \omterm{def:rs:point-in-time-data-and-the-biases:bitemporal}{knowledge times}, store splits as events with theirs (chapter~4), and rescale at
query time to the share basis of the decision date, the basis of the price the number will be divided by.

\begin{omfigure}
\begin{tikzpicture}
\begin{semilogyaxis}[width=11cm,height=5.6cm,xlabel={quarter end},ylabel={TTM diluted EPS (\$)},tick label style={font=\small},
  label style={font=\small},xmin=2011.5,xmax=2026.7,ymin=2,ymax=80,ytick={2,5,10,20,50},yticklabels={2,5,10,20,50},
  grid=major,grid style={gray!25},
  xticklabel style={/pgf/number format/1000 sep=},legend pos=north east,legend style={font=\footnotesize},table/col sep=comma]
\addplot[omThm,thick,mark=*,mark size=1.2pt] table[x=year,y=first]{figdata/research/11-fundamental-analyst-and-event-features/apple_ttm.csv};
\addplot[omDef,thick,dashed] table[x=year,y=latest]{figdata/research/11-fundamental-analyst-and-event-features/apple_ttm.csv};
\draw[black!50] (axis cs:2014.43,2) -- (axis cs:2014.43,80) node[pos=0.06,right,font=\footnotesize]{7:1};
\draw[black!50] (axis cs:2020.66,2) -- (axis cs:2020.66,80) node[pos=0.06,right,font=\footnotesize]{4:1};
\legend{from first-reported quarters,from today's values}
\end{semilogyaxis}
\end{tikzpicture}

\omcaption{Apple's trailing-twelve-month diluted EPS at each quarter's end, summed from each quarter's first-reported value and from
its latest value in EDGAR, with the two stock splits. Neither is a usable series without split adjustment at the decision date. Data: SEC
EDGAR company facts.}\label{fig:rs:fundamental-analyst-and-event-features:ttm}
\end{omfigure}

\section{Ratios as features}

\begin{definition}[Accounting ratio]\label{def:rs:fundamental-analyst-and-event-features:ratio}
An \emph{accounting ratio}\index{accounting ratio} divides one accounting quantity by another or by a market quantity: book-to-price,
earnings yield (TTM earnings over price), sales-to-price, return on equity, leverage, asset growth.
\end{definition}

A ratio has two \omterm{def:rs:point-in-time-data-and-the-biases:bitemporal}{knowledge times} and must respect both. The accounting side is known at its filing (or release); the market side
(price, market capitalisation) is known at its own time, normally the decision time. The pairings that fail are regular: a TTM earnings
figure divided by a price on a different share basis; a book value keyed on its period end; a ratio whose denominator can change sign
(earnings, book equity), which needs a separate treatment for negative values rather than a rank that puts the most negative earnings
yield next to the most expensive stock. Book-to-price in \texttt{firm.synthmkt} shows how much the key date matters for a slow variable:
its mean rank IC over the next 21 days is 0.029 (standard error 0.038) keyed on the period end and $-0.005$ (0.023) keyed on the filing
day. The difference is noise: each quarterly IC is dominated by the value factor's realised return over its window, and moving the window
a few weeks changes little in expectation because nothing in the simulated market happens on the filing day. The damage done by a wrong key
date is proportional to how much the number's release moves the price, and a book value's release does not move it much.

\section{Surprises and revisions}

\begin{definition}[Earnings surprise, standardised unexpected earnings]\label{def:rs:fundamental-analyst-and-event-features:sue}
An \emph{earnings surprise}\index{earnings surprise} is the reported earnings of a period minus their expected value: an analysts' consensus
forecast (Book~2, chapter~31) or a time-series model. \emph{Standardised unexpected earnings}\index{standardised unexpected earnings} (SUE)
divide the surprise of a seasonal random walk, $x_q - x_{q-4}$, by the standard deviation of that difference over the previous eight
quarters.
\end{definition}

Foster, Olsen and Shevlin (1984) documented that systematic post-announcement drifts in returns are associated with the sign and magnitude of
unexpected earnings, and that for expectation models based on the time series of reported quarterly earnings the forecast error and firm
size explained 81\% and 61\% of the variation in the drifts. The release is the event; the drift after it is the feature's opportunity; a
key date before the release captures the release itself. In \texttt{firm.synthmkt} each quarter's earnings are announced a random number of trading days after its end (31 on average, about six weeks: the head start of a period-end key), with a jump in the direction of the true
surprise and a smaller drift afterwards. The surprise keyed on its announcement has a mean rank IC with the next 21 days' market-adjusted
return of 0.020 (standard error 0.006): the planted drift. Keyed on the period end it is 0.167, eight times larger, because a third of the
windows contain the announcement's jump (\cref{fig:rs:fundamental-analyst-and-event-features:ic}).

\begin{omfigure}
\begin{tikzpicture}
\begin{axis}[width=10cm,height=5.4cm,ybar,bar width=16pt,ylabel={mean rank IC, next 21 days},tick label style={font=\small},
  label style={font=\small},xtick={0,1,2,3},xticklabels={{SUE\\known},{SUE\\period end},{B/P\\known},{B/P\\period end}},
  xticklabel style={align=center},
  ymin=-0.1,ymax=0.2,enlarge x limits=0.18,grid=major,grid style={gray!25},table/col sep=comma]
\addplot[fill=omDef!55,draw=omDef,error bars/.cd,y dir=both,y explicit]
  table[x=k,y=ic,y error=se2]{figdata/research/11-fundamental-analyst-and-event-features/ic.csv};
\end{axis}
\end{tikzpicture}

\omcaption{Mean rank IC over 39 quarterly cross-sections of \texttt{firm.synthmkt} (\omterm{def:rs:market-data-for-research:bar}{bars}: two standard errors) of the true \omterm{def:rs:fundamental-analyst-and-event-features:sue}{earnings surprise}
and of book-to-price, keyed on the day each became known (announcement, filing) and on the period's end. The period-end key multiplies the
surprise's IC by eight and leaves book-to-price within its noise.}\label{fig:rs:fundamental-analyst-and-event-features:ic}
\end{omfigure}

\begin{definition}[Analyst revision, forecast dispersion]\label{def:rs:fundamental-analyst-and-event-features:revision}
An \emph{analyst revision}\index{analyst revision} is a change in an analyst's forecast or recommendation for a security; as a feature, the
change in the consensus forecast over a window, scaled by price or by the forecasts' dispersion. \emph{Forecast
dispersion}\index{forecast dispersion} is the cross-analyst standard deviation of the current forecasts divided by the absolute value of their
mean.
\end{definition}

Womack (1996) found that new buy and sell recommendations by analysts at major US brokerages move prices at once, although few coincide
with new public news, and that the reactions are incomplete: the mean drift after a buy recommendation was modest (+2.4\%) and short-lived,
after a sell recommendation larger ($-9.1\%$) and six months long. Diether, Malloy and Scherbina (2002) found that stocks with higher dispersion
of analysts' earnings forecasts earn lower future returns than otherwise similar stocks, most so in small stocks and past losers, consistent
with prices reflecting the optimists' view when the most pessimistic investors do not trade. Both are point-in-time problems first: a
forecast database must keep each forecast with the date the analyst made it and the date the vendor recorded it, and a recommendation that
is later withdrawn or corrected must stay in history as it was. Coverage is its own bias: analysts cover the firms they expect to do well
and drop the others, so the set of covered firms changes with the outcome.

\section{Event calendars}

\begin{definition}[Event calendar, announcement window]\label{def:rs:fundamental-analyst-and-event-features:calendar}
An \emph{event calendar}\index{event calendar} lists scheduled corporate events (earnings releases, dividend dates, shareholder meetings,
index changes) with the date each was announced and the date it occurs. An \emph{announcement window}\index{announcement window} is a fixed
range of trading days around an event, such as $[-1, +1]$, used to measure its reaction or to exclude it.
\end{definition}

A calendar is itself \omterm{def:rs:point-in-time-data-and-the-biases:lookahead}{point-in-time data}: a scheduled date is announced before the event, may move, and is stored with the time it
was announced; a date that moves is itself worth testing as a feature. The ten companies' release lags are regular: Home Depot released
16 days after the end of every one of its first three fiscal quarters since 2013, Exxon Mobil 30 days give or take two, while Apple's varies more (a standard deviation of five days). A predicted date is therefore a usable feature before the
confirmed date exists, and \emph{days to the next event} is a conditioning variable every research system needs: it
switches features on (a surprise after the release), off (a volatility estimate across a release), and sizes risk around the known unknown.
\omterm{def:rs:fundamental-analyst-and-event-features:calendar}{Announcement windows} serve the other way round: a reversal or volatility feature estimated through an earnings day measures the announcement,
not the effect it was meant to measure, so research code excludes the window or models it separately.

\section{Predictor cards}

\begin{predictorcard}{Standardised unexpected earnings}\label{pred:rs:fundamental-analyst-and-event-features:sue}
\sfield{Definition} $(x_q - x_{q-4})/\sigma_8$, the seasonal random-walk surprise of quarterly EPS over the standard deviation of the last
eight such differences; or the surprise against the consensus, over the dispersion.
\sfield{Inputs and timestamps} Quarterly EPS as first released, keyed on the release (the 8-K) or, from filings only, on the filing; split
factors as known at the decision date.
\sfield{Rationale} Post-announcement drift (Foster, Olsen and Shevlin, 1984).
\sfield{Horizon and half-life} Weeks to a quarter after the release; on \texttt{firm.synthmkt}, IC 0.020 over 21 days with the correct key.
\sfield{Normalisation} Rank across the stocks that announced in the window; neutralise size.
\sfield{Failure modes} A period-end key (IC 0.167 on the simulated market: all of it the announcement); a surprise measured against a
restated number; a stale consensus.
\sfield{Sources} As cited; \texttt{firm.fundpit.sue}; \texttt{rs\_fundamentals.ic\_inflation}.
\end{predictorcard}

\begin{predictorcard}{Book-to-price, as known}\label{pred:rs:fundamental-analyst-and-event-features:bp}
\sfield{Definition} The latest book equity known at the decision time, over the market capitalisation at the decision time.
\sfield{Inputs and timestamps} Book equity keyed on its filing; shares and price at the decision date, on one share basis.
\sfield{Rationale} The value premium (Book~8, chapter~6).
\sfield{Horizon and half-life} Months to years; the book changes quarterly, the price daily.
\sfield{Normalisation} Rank within industry; negative book equity handled apart.
\sfield{Failure modes} A period-end key (little damage for a slow ratio: on the simulated market the difference is within noise);
mismatched share bases; restated book values used before their \omterm{def:rs:point-in-time-data-and-the-biases:vintage}{restatement}.
\sfield{Sources} \texttt{rs\_fundamentals.ic\_inflation}.
\end{predictorcard}

\section{Tutorial: the five-week head start}\label{tut:rs:fundamental-analyst-and-event-features:tut}

\begin{tutorial}
\textbf{Goal.} Build the dates of real earnings numbers from EDGAR, store them point-in-time, derive quarters and TTM values as known at
any date, and measure on a synthetic market what a period-end key does to a surprise and a ratio. \textbf{End state:}
\cref{fig:rs:fundamental-analyst-and-event-features:calendar,fig:rs:fundamental-analyst-and-event-features:ttm,fig:rs:fundamental-analyst-and-event-features:ic};
the \omterm{def:rs:point-in-time-data-and-the-biases:vintage}{restatement} table.
\begin{tutsteps}
\item \textbf{Fetch} with \texttt{rs\_fetch\_edgar.py}: submissions (8-K item 2.02, 10-Q, 10-K) and company facts (diluted EPS with every filing
that reported it), with a descriptive User-Agent and a pause between requests.
\item \textbf{Store} each fact under its duration, valid at the period end, known at the filing.
\omcode{code/firm/fundpit/firm_fundpit.py}{42}{54}{Facts into a point-in-time store, by duration.}\label{lst:rs:fundamental-analyst-and-event-features:load}
\item \textbf{Quarters as known}, the fourth derived from the year and the nine months.
\omcode{code/firm/fundpit/firm_fundpit.py}{57}{72}{Quarterly values as known at a date.}\label{lst:rs:fundamental-analyst-and-event-features:quarterly}
\item \textbf{Run} \texttt{calendar\_summary()}, \texttt{restatements()}, \texttt{apple\_ttm()}, \texttt{ic\_inflation()} and
\texttt{fig\_fundamentals.py}.
\end{tutsteps}
\textbf{What to change next.} Key the surprise on the 10-Q filing instead of the release and measure what the week between them costs;
rescale Apple's first-reported quarters to the share basis of each decision date and rebuild a clean TTM series.
\end{tutorial}

\section{Build: the fundamental feature store}\label{bld:rs:fundamental-analyst-and-event-features:fundpit}

\begin{build}
\bfield{Purpose} Accounting features for the miniature firm's research that cannot leak: every value carries its period and its \omterm{def:rs:point-in-time-data-and-the-biases:bitemporal}{knowledge
time}, and every derived value (quarters, TTM, SUE, ratios) is computed from what was known at the decision date.
\bfield{Interface} \texttt{duration\_kind}, \texttt{load(facts, field)}, \texttt{quarterly(store, entity, field, known)}, \texttt{ttm},
\texttt{sue(values, n)}, \texttt{restated}, \texttt{split\_ratio}, \texttt{revision}, \texttt{dispersion}, \texttt{days\_to\_event},
\texttt{in\_window}; storage is \texttt{firm.pit.Store}.
\bfield{Rules} \omterm{def:rs:point-in-time-data-and-the-biases:bitemporal}{Valid time} is the period end, \omterm{def:rs:point-in-time-data-and-the-biases:bitemporal}{knowledge time} the filing (or the release when it is stored); a fourth quarter exists only when
the year and the nine months are both known; a per-share value is never compared across share bases.
\bfield{Acceptance tests} \texttt{code/firm/fundpit/tests/}: durations of a hand-built year, the fourth quarter (\$4.80 minus \$3.30), a quarter
unknown before its filing, TTM, the later view after a split; \omterm{def:rs:point-in-time-data-and-the-biases:vintage}{restatement} listing and split ratios to the cent; SUE of a seasonal jump and of
a random walk; revision, dispersion, days to the next event and windows by hand.
\bfield{Stretch} Release dates from 8-K filings as the \omterm{def:rs:point-in-time-data-and-the-biases:bitemporal}{knowledge time} of each quarter; a consensus store with analyst and vendor timestamps;
the split events of \texttt{firm.secmaster} applied at query time.
\end{build}

\begin{omsources}
\item US Securities and Exchange Commission, Form 10-Q and Form 10-K, general instructions; EDGAR application programming interfaces
(data.sec.gov).
\item Apple Inc., Form 10-K for fiscal 2014 and fiscal 2020; Microsoft Corporation, Form 10-K for fiscal 2018 (SEC EDGAR).
\item G.~Foster, C.~Olsen and T.~Shevlin, ``Earnings releases, anomalies, and the behavior of security returns'', \emph{The Accounting Review}
59(4), 1984.
\item K.~L.~Womack, ``Do brokerage analysts' recommendations have investment value?'', \emph{Journal of Finance} 51(1), 1996.
\item K.~B.~Diether, C.~J.~Malloy and A.~Scherbina, ``Differences of opinion and the cross section of stock returns'', \emph{Journal of
Finance} 57(5), 2002.
\end{omsources}

\section{Exercises}

\begin{exercise}[$\star$]\label{exo:rs:fundamental-analyst-and-event-features:1}
A company's fiscal year ends on 30 June. It reports nine-month diluted EPS of \$3.30 at 31 March and annual EPS of \$4.80. What is the fourth
quarter's EPS, and when is it first known if the 10-K is filed 55 days after the year's end?
\end{exercise}

\begin{exercise}[$\star$]\label{exo:rs:fundamental-analyst-and-event-features:2}
An accelerated filer's quarter ends on 30 September. By which date must its 10-Q be filed? And a large accelerated filer's 10-K for a year
ending 31 December (not a leap year)?
\end{exercise}

\begin{exercise}[$\star$]\label{exo:rs:fundamental-analyst-and-event-features:3}
A quarter's EPS was first reported as \$1.42 and later as \$0.355. What happened, and how should the store keep both?
\end{exercise}

\begin{exercise}[$\star\star$]\label{exo:rs:fundamental-analyst-and-event-features:4}
Quarterly EPS over twelve quarters are 1.00, 1.10, 0.90, 1.20, 1.04, 1.12, 0.96, 1.22, 1.06, 1.16, 0.94, 1.40. Compute the seasonal differences and the SUE of the last quarter with $n = 6$ (the six previous differences, sample standard deviation).
\end{exercise}

\begin{exercise}[$\star\star$]\label{exo:rs:fundamental-analyst-and-event-features:5}
If announcements fall uniformly between 0 and 62 trading days after a quarter's end, what share of 21-day windows starting the day after the
quarter's end contain the announcement? Relate it to the eightfold IC inflation.
\end{exercise}

\begin{exercise}[$\star\star$]\label{exo:rs:fundamental-analyst-and-event-features:6}
Why can book-to-price keyed on the period end show little inflation while the \omterm{def:rs:fundamental-analyst-and-event-features:sue}{earnings surprise} shows a lot? What would make book-to-price
inflate too?
\end{exercise}

\begin{exercise}[$\star\star\star$]\label{exo:rs:fundamental-analyst-and-event-features:7}
\emph{Coding.} From \texttt{data/research/edgar\_filings.csv}, compute for each company the share of quarters in which the 10-Q was filed on
the release day, and the median days from release to filing for the others.
\end{exercise}

\begin{exercise}[$\star\star\star$]\label{exo:rs:fundamental-analyst-and-event-features:8}
\emph{Find the flaw.} ``We use the vendor's historical fundamentals table, which is updated whenever a company restates, so our earnings
history is always the most accurate available.''
\end{exercise}

\section{Problem: The Five-Week Head Start}

\begin{problem}[{Weekend problem --- a key date, measured}]\label{pb:rs:fundamental-analyst-and-event-features:1}
Ten companies' EDGAR data since 2013, and \texttt{firm.synthmkt} with its earnings announcements, filings and book values.

\medskip\noindent\textbf{Part I --- The dates.}
\begin{enumerate}
\item How many period ends have the ten companies since 2013, and how many are followed by a release within 90 days?
\item What are the median days from a quarter's end to the release and to the 10-Q, and for a fiscal year to the release and the 10-K?
\item What are the largest 10-Q and 10-K lags, and how do they compare with the deadlines?
\item How often is the report filed on the release day, and what is the median gap?
\end{enumerate}

\medskip\noindent\textbf{Part II --- The values.}
\begin{enumerate}[resume]
\item How many quarterly values changed after their first filing, and how many by an exact ratio?
\item What were Apple's two splits, and what did they do to its per-share history?
\item What happened to Microsoft's fiscal 2017 EPS, and why?
\item What TTM does Apple's fiscal 2014 give from first-reported quarters, and what annual EPS did Apple report?
\item How should a store keep split-affected values?
\end{enumerate}

\medskip\noindent\textbf{Part III --- The damage.}
\begin{enumerate}[resume]
\item What is the average head start of a period-end key in \texttt{firm.synthmkt}?
\item What are the surprise's IC keyed on the announcement and on the period end, with their standard errors?
\item What share of the windows contain the announcement?
\item What are book-to-price's ICs under the two keys, with standard errors?
\item Why is the book-to-price difference not significant?
\item State the \emph{named result}: the IC inflation of a surprise and a value signal when period-end dating replaces known-date dating.
\end{enumerate}

\medskip\noindent\textbf{Part IV --- Beyond.}
\begin{enumerate}[resume]
\item Key the surprise on the filing instead of the release: what do you expect to lose?
\item Which is worse for a backtest, a period-end key or a latest-value (restated) history, for a SUE feature?
\item How would you detect a period-end key in someone else's backtest?
\item What should a fundamentals vendor deliver for research use?
\item In one sentence: what is the \omterm{def:rs:point-in-time-data-and-the-biases:bitemporal}{knowledge time} of an accounting number?
\end{enumerate}
\end{problem}

\section{Interview questions}

\begin{interviewq}[$\star$ \iqroles{researcher}]\label{iq:rs:fundamental-analyst-and-event-features:1}
What is point-in-time fundamental data, and why does it matter?
\end{interviewq}

\begin{interviewq}[$\star\star$ \iqroles{researcher}]\label{iq:rs:fundamental-analyst-and-event-features:2}
How would you compute trailing-twelve-month earnings for a company that reports only year-to-date figures?
\end{interviewq}

\begin{interviewq}[$\star\star$ \iqroles{researcher, trader}]\label{iq:rs:fundamental-analyst-and-event-features:3}
What is post-earnings-announcement drift, and how would you build a feature for it?
\end{interviewq}

\begin{interviewq}[$\star\star$ \iqroles{researcher}]\label{iq:rs:fundamental-analyst-and-event-features:4}
Your value signal's backtest improves by a third when you switch from the filing date to the period end as the key. Which is right, and
why might the improvement be small or large?
\end{interviewq}

\begin{interviewq}[$\star\star$ \iqroles{researcher}]\label{iq:rs:fundamental-analyst-and-event-features:5}
What biases affect analyst forecast data?
\end{interviewq}

\begin{interviewq}[$\star\star\star$ \iqroles{researcher, developer}]\label{iq:rs:fundamental-analyst-and-event-features:6}
Design the storage of per-share accounting values so that any query returns what was known at a date, on the share basis of that date.
\end{interviewq}
